% TEMPLATE-MILC-v3.tex - plantilla maestra UNICA de las guias MILC v3.
%
% La usan las tres familias de guias, cada una con su driver:
%   · Grados 6-11  -> scripts/build-guias-g11.py
%   · Bebras       -> scripts/build-guias-bebras.py
%   · Semillero    -> scripts/build-guias-semillero.py
%
% Antes habia tres copias casi identicas (~400 lineas iguales, 6 distintas):
% cada mejora habia que aplicarla tres veces, y en la practica no se hacia
% (el motor de recursos vivia solo en dos de ellas). Lo que varia entre
% familias esta parametrizado en 6 placeholders que cada driver llena
% (se nombran aqui SIN su sintaxis real: el builder reemplaza por texto plano,
% tambien dentro de los comentarios, y un valor multilinea rompe el preambulo):
%
%   HEADER_IZQ        encabezado izquierdo de pagina
%   HEADER_DER        encabezado derecho de pagina
%   PORTADA_ETIQUETA  rotulo sobre el numero grande (GRADO/MOMENTO/MODULO)
%   PORTADA_SERIE     linea de serie de la portada
%   PORTADA_ID        identificador de la guia en la portada
%   PORTADA_CIRCULO   el circulito de grado (°), o vacio si no aplica
%
% El resto de claves las documenta content/guias/_SCHEMA.md. El script valida
% que no queden placeholders sin reemplazar antes de escribir el archivo final.

\documentclass[11pt,letterpaper]{article}
\usepackage[margin=1.75cm]{geometry}
\usepackage[table]{xcolor}
\usepackage{fontspec}
\usepackage[spanish]{babel}
\usepackage{tikz}
% Librerías TikZ para los diagramas de las guías (bloque `recursos:`):
%  · babel  → babel en español vuelve activos caracteres como «>», y TikZ los usa
%             en sus opciones (>=stealth, ->). Sin esta librería, cualquier
%             diagrama con flechas revienta con «\language@active@arg> has an
%             extra }». Debe ir ANTES de que se dibuje nada.
%  · positioning        → sintaxis «below left=2pt and 3pt».
%  · decorations.pathreplacing → llaves ({brace}) para señalar rangos.
\usetikzlibrary{babel,positioning,decorations.pathreplacing}
\usepackage{graphicx}
% Circuitos: el trabajo de electrónica se hace hoy con micro:bit y sus sensores
% integrados (MakeCode, sin cableado), así que no se necesitan esquemáticos.
% Si algún día entran montajes con protoboard, basta descomentar esta línea:
% los diagramas .tex/.tikz declarados en `recursos:` ya se incrustan solos.
% \usepackage{circuitikz}
\usepackage[most]{tcolorbox}
\usepackage{tabularx}
\usepackage{array}
\usepackage{fancyhdr}
\usepackage{titlesec}
\usepackage{ragged2e}
\usepackage{enumitem}
\usepackage{microtype}

% Helvetica Neue no trae la flecha U+2192, asi que cada «→» escrito literalmente
% en el YAML se compilaba a NADA, en silencio (462 flechas en 61 guias: rutas del
% tipo «paleta → tipografia → lockup» salian rotas). Mapeamos el caracter a su
% version matematica, que si tiene glifo. Asi el autor puede seguir escribiendo
% «→» comodamente en el YAML.
\usepackage{newunicodechar}
\newunicodechar{→}{\ensuremath{\rightarrow}}

\setmainfont{Helvetica Neue}
\setsansfont{Helvetica Neue}
\setlength{\parindent}{0pt}
\setlength{\parskip}{6pt}
\hyphenpenalty=200
\exhyphenpenalty=200
\pretolerance=400
\tolerance=2000
\setlength{\emergencystretch}{1em}
\renewcommand{\arraystretch}{1.25}
\setlist{leftmargin=5mm,itemsep=1mm,topsep=1mm}
\newcolumntype{Y}{>{\RaggedRight\arraybackslash}X}

\definecolor{milcMagenta}{HTML}{E6007E}
\definecolor{milcVino}{HTML}{5A0038}
\definecolor{milcTurquesa}{HTML}{0093A5}
\definecolor{milcVerde}{HTML}{58A923}
\definecolor{milcMostaza}{HTML}{E5B400}
\definecolor{milcOcre}{HTML}{B86600}
\definecolor{milcNegro}{HTML}{241816}
\definecolor{milcGris}{HTML}{F7F7F4}
\definecolor{milcLinea}{HTML}{E8E2DD}
\definecolor{milcGradePrimary}{HTML}{FF2D87}
\definecolor{milcGradeDark}{HTML}{9F1B56}
\definecolor{milcGradeSoft}{HTML}{FFE0EE}
\definecolor{milcGradeAccent}{HTML}{CFE7FF}
\definecolor{milcGradeText}{HTML}{FFFFFF}
\color{milcNegro}

\pagestyle{fancy}
\fancyhf{}
\lhead{\small Tecnología e Informática · Grado 10}
\rhead{\small Guía 9 · MILC v3}
\cfoot{\small\thepage}
\renewcommand{\headrulewidth}{0pt}

\newtcolorbox{titlebox}[1]{
  enhanced,colback=#1,colframe=#1,arc=7mm,boxrule=0pt,
  left=7mm,right=7mm,top=3mm,bottom=3mm,
  before skip=8pt,after skip=8pt,
  fontupper=\sffamily\bfseries\Large\color{white}
}

\newtcolorbox{softbox}[3]{
  enhanced,colback=#2,colframe=#1,borderline west={4pt}{0pt}{#1},
  arc=2mm,boxrule=.4pt,left=4mm,right=4mm,top=3mm,bottom=3mm,
  before skip=6pt,after skip=8pt,title={#3},coltitle=white,
  colbacktitle=#1,fonttitle=\sffamily\bfseries,fontupper=\RaggedRight,
  attach boxed title to top left={xshift=3mm,yshift=-2mm},
  boxed title style={arc=2mm, boxrule=0pt}
}

\newtcolorbox{drawbox}[2]{
  enhanced,colback=white,colframe=#1,arc=4mm,boxrule=1pt,
  left=5mm,right=5mm,top=4mm,bottom=4mm,before skip=8pt,after skip=8pt,
  title={#2},coltitle=white,colbacktitle=#1,fonttitle=\sffamily\bfseries,
  fontupper=\RaggedRight,
  attach boxed title to top left={xshift=4mm,yshift=-2mm},
  boxed title style={arc=3mm, boxrule=0pt}
}

\newtcolorbox{infoband}[1]{
  enhanced,colback=#1!8,colframe=#1!50,arc=2mm,boxrule=.5pt,
  left=4mm,right=4mm,top=2mm,bottom=2mm,before skip=4pt,after skip=4pt,
  fontupper=\small\RaggedRight
}

% Puente narrativo entre secciones (contrato editorial Conéctate).
% Da continuidad: "Acabas de X. Ahora vas a Y porque Z."
% Visualmente: banda izquierda en color del grado, texto en itálica pequeña.
\newtcolorbox{puentebox}{
  enhanced,colback=milcGradeSoft,colframe=milcGradeDark,
  borderline west={3pt}{0pt}{milcGradeDark},
  arc=0mm,boxrule=0pt,left=5mm,right=4mm,top=2.5mm,bottom=2.5mm,
  before skip=4pt,after skip=8pt,
  fontupper=\itshape\small\color{milcNegro}\RaggedRight
}
% La flecha va en modo matematico a proposito: Helvetica Neue NO tiene el glifo
% U+2192, asi que \textrightarrow se compilaba a NADA («Missing character: There
% is no → in font Helvetica Neue») y el puente salia sin flecha. En modo
% matematico la flecha viene de la fuente matematica y si se dibuja.
\newcommand{\puente}[1]{%
  \begin{puentebox}%
    \textcolor{milcGradeDark}{$\rightarrow$}\hspace{2mm}#1%
  \end{puentebox}%
}

\newcommand{\linead}[1]{\noindent\color{milcLinea}\rule{#1}{0.5pt}\color{milcNegro}}
\newcommand{\respuesta}[1]{\par\vspace{2mm}\foreach \n in {1,...,#1}{\noindent\color{milcLinea}\rule{\linewidth}{0.5pt}\par\vspace{5mm}}\color{milcNegro}}
\newcommand{\checkbox}{\textcolor{milcGradeDark}{\fbox{\phantom{x}}}\hspace{5pt}}

\newcommand{\phasecard}[4]{
\begin{tcolorbox}[enhanced,colback=#3,colframe=#2,arc=3mm,boxrule=.5pt,height=3.05cm,valign=center,halign=center,left=1.5mm,right=1.5mm,top=2mm,bottom=2mm]
{\sffamily\bfseries\fontsize{22}{24}\selectfont\textcolor{#2}{#1}}\par
{\sffamily\bfseries\small #4}
\end{tcolorbox}
}

% ── Recursos visuales de la guía ──────────────────────────────────────
% Declarados en el bloque `recursos:` del YAML y emitidos por el builder.
% \guiaFigura   → imagen rasterizada o PDF (foto, carta celeste, montaje).
% \guiaDiagrama → fuente TikZ/circuitikz (.tex) incrustada como vector:
%                 es la vía para circuitos y esquemáticos editables.
% #1 = ruta relativa al .tex · #2 = pie (puede ir vacío)
\newcommand{\guiaFigura}[2]{%
\begin{center}
\includegraphics[width=0.88\linewidth,height=0.38\textheight,keepaspectratio]{#1}\\[1.5mm]
{\footnotesize\itshape\textcolor{milcNegro!75}{#2}}
\end{center}
}

\newcommand{\guiaDiagrama}[2]{%
\begin{center}
\input{#1}\\[1.5mm]
{\footnotesize\itshape\textcolor{milcNegro!75}{#2}}
\end{center}
}

\begin{document}
\thispagestyle{empty}
\begin{tikzpicture}[remember picture,overlay]
  \fill[white] (current page.south west) rectangle (current page.north east);
  \path[fill=milcGradePrimary,rounded corners=16mm]
    ([xshift=.55cm,yshift=-.55cm]current page.north west) rectangle
    ([xshift=-.55cm,yshift=3.65cm]current page.south east);

  \node[anchor=north west,text=milcGradeText] at ([xshift=2.0cm,yshift=-1.85cm]current page.north west)
    {\sffamily\bfseries\fontsize{14}{16}\selectfont GRADO};
  \node[anchor=north west,text=milcGradeText,opacity=.82] at ([xshift=2.0cm,yshift=-2.75cm]current page.north west)
    {\sffamily\bfseries\fontsize{9}{11}\selectfont SERIE GUÍAS · TECNOLOGÍA E INFORMÁTICA};

  \begin{scope}[shift={([xshift=-3.05cm,yshift=-2.55cm]current page.north east)}]
    \fill[white,opacity=.80] (-.62,-.52) rectangle (.62,.52);
    \draw[milcGradeText,line width=1pt] (-.62,-.52) rectangle (.62,.52);
    \fill[milcGradeDark] (-.42,-.38) rectangle (-.22,.06);
    \fill[milcGradeAccent] (-.10,-.38) rectangle (.10,.24);
    \fill[milcGradePrimary!60!black] (.22,-.38) rectangle (.42,.38);
  \end{scope}

  \node[anchor=west,text=milcGradeText] at ([xshift=1.9cm,yshift=-12.85cm]current page.north west)
    {\sffamily\bfseries\fontsize{124}{124}\selectfont 10};
  % El circulito de grado (°) solo aplica a las guias de grado («11°»); en
  % Bebras y Semillero el numero grande es un momento/modulo, no un grado.
  \draw[milcGradeText,line width=1.7pt]
    ([xshift=4.52cm,yshift=-11.68cm]current page.north west) circle (.13cm);

  \node[anchor=west,text=milcGradeText,text width=8.7cm,align=left] at ([xshift=10.35cm,yshift=-11.6cm]current page.north west)
    {
     {\sffamily\bfseries\fontsize{19}{22}\selectfont Producción final ---\\80 páginas compiladas,\\revisadas y firmadas}\\[4mm]
     {\sffamily\bfseries\fontsize{23}{26}\selectfont Guía 9-10-TIC}\\[2mm]
     {\sffamily\fontsize{13}{16}\selectfont Período 1 · Escritura abierta con IA}\\[5mm]
     {\sffamily\bfseries\fontsize{11}{13}\selectfont Institución Educativa Sor María Juliana}\\[1mm]
     {\sffamily\fontsize{10}{12}\selectfont Autor: Dr. Álvaro Cárdenas Orozco}
    };

  \path[fill=milcGradeSoft,rounded corners=7mm]
    ([xshift=1.35cm,yshift=.85cm]current page.south west) rectangle
    ([xshift=-1.35cm,yshift=4.25cm]current page.south east);
  \draw[milcGradeDark,line width=.65pt,rounded corners=7mm]
    ([xshift=1.35cm,yshift=.85cm]current page.south west) rectangle
    ([xshift=-1.35cm,yshift=4.25cm]current page.south east);
  \node[anchor=center,text=milcNegro,text width=17.0cm,align=center] at ([yshift=2.55cm]current page.south)
    {\sffamily\fontsize{10}{12}\selectfont Metodología MILC · Apertura ancestral · Pensamiento computacional · Triángulo Dussel-Estoicismo-Floridi\\
    \textcolor{milcGradeDark}{\bfseries Producto:} Libro completo de al menos 80 páginas en formato PDF compilado integrando portada + tabla de contenido + 8-12 capítulos diagramados + página de créditos con declaración de derechos y uso de IA + carta del editor (tú) de 1 página firmada con nombre y fecha + revisión de calidad pasando los 12 puntos del checklist + bitácora del proceso completo del periodo documentando los 9 pasos recorridos hasta hoy};
\end{tikzpicture}
\mbox{}
\newpage

\section*{Apertura: Producción final del libro --- 80 páginas compiladas, revisadas y firmadas}

\begin{softbox}{milcGradeDark}{milcGradeSoft}{Saber ancestral}
En los talleres antiguos de carpintería, cerámica y orfebrería, hubo durante siglos un momento ritual que separaba al aprendiz del oficial: \textbf{la firma de la pieza maestra}. Después de un año de trabajo bajo el maestro, el aprendiz producía una pieza final que demostraba el oficio aprendido: una mesa con todas las uniones, un jarrón con todas las técnicas, un anillo con todos los procesos. Esa pieza no se entregaba cualquier día. Había un acto formal en el taller con el maestro y los oficiales viejos sentados alrededor. El aprendiz traía la pieza, la ponía sobre la mesa, la explicaba: \emph{``aquí usé esta técnica, esto lo aprendí del maestro, esto lo decidí yo, esto aún no domino del todo''}. Los mayores hacían preguntas. Si la defensa era sólida, el aprendiz tomaba un punzón y \textbf{firmaba la pieza}: grabar su nombre o marca en la base. Esa firma significaba 3 cosas al mismo tiempo: \textbf{declarar autoría, asumir responsabilidad, certificar el oficio}. La pieza con firma era certificación de competencia profesional. El libro de 80 páginas que entregas hoy es exactamente esa pieza maestra: producto final del aprendizaje del periodo, firmado con tu nombre, defendible ante cualquier audiencia.
\end{softbox}

\begin{softbox}{milcTurquesa}{milcGris}{Saber contemporáneo}
La \textbf{producción final del libro} integra todo lo recorrido en las 8 sesiones anteriores en un \textbf{PDF compilado de al menos 80 páginas} listo para circular. Sus componentes obligatorios son: (1) \textbf{Portada} (sesión 7): con título, autor, imagen. (2) \textbf{Página de créditos}: año, autor, licencia de derechos (sesión 6), declaración de uso de IA. (3) \textbf{Tabla de contenido autogenerada} (sesión 8). (4) \textbf{8-12 capítulos diagramados} (sesiones 4-5-8): cada uno con texto V3, diagramación coherente, numeración correcta. (5) \textbf{Carta del editor} (1 página firmada): tu palabra como editor responsable del libro. (6) \textbf{Contraportada} opcional con sinopsis. La regla profesional: \textbf{el libro no se entrega antes de pasar checklist de revisión}. El checklist típico tiene 12 puntos que cubren contenido, forma, derechos, voz propia. La \textbf{carta del editor} es pieza nueva y esencial: declaración personal sobre el proceso del libro, el uso de IA, lo que se logró y lo que aún se considera mejorable. Firmada con nombre y fecha, es tu sello profesional sobre la pieza maestra.
\end{softbox}

\begin{softbox}{milcMostaza}{milcGris}{Pregunta-puente}
\textit{¿Qué sabía el aprendiz al grabar su nombre en la pieza maestra ante el maestro y los oficiales viejos, que el estudiante novato olvida cuando entrega un PDF sin haber revisado ni firmado? ¿Y por qué la carta del editor es la pieza más profesional del libro, aunque sea solo 1 página?}
\end{softbox}

\begin{softbox}{milcVerde}{milcGris}{Saber y saber hacer}
Al terminar podrás: (1) \textbf{analizar} el material producido en las 8 sesiones anteriores y verificar que esté listo para integración; (2) \textbf{evaluar} el libro completo con el checklist de 12 puntos antes de declararlo final; (3) \textbf{crear} el PDF compilado de 80+ páginas con todos los componentes y la carta del editor firmada; (4) \textbf{explicar} en la carta del editor el proceso completo del libro con honestidad sobre uso de IA y decisiones tomadas.
\end{softbox}



\small
\begin{tabularx}{\textwidth}{>{\bfseries}p{3.0cm}Y}
\rowcolor{milcGradePrimary!12}Autor & Dr. Álvaro Cárdenas Orozco\\
DBA & Producción editorial mediada por IA con criterio ético y técnico (MEN, T\&I)\\
Referentes & Ley 23 de 1982 · Floridi (ética de la IA generativa) · Dussel (autoría con dignidad)\\
Producto final & Libro completo de al menos 80 páginas en formato PDF compilado integrando portada + tabla de contenido + 8-12 capítulos diagramados + página de créditos con declaración de derechos y uso de IA + carta del editor (tú) de 1 página firmada con nombre y fecha + revisión de calidad pasando los 12 puntos del checklist + bitácora del proceso completo del periodo documentando los 9 pasos recorridos hasta hoy\\
\end{tabularx}
\normalsize

\puente{Antes de empezar, mira el viaje completo. Has recorrido 8 sesiones: editor, concepción, prompting, iteración, estructura, derechos, portada, diagramación. Hoy es el \textbf{día de la firma}: integras todo en el libro compilado. La sesión 10 será sustentación pública con feria del libro escolar.}

\begin{titlebox}{milcGradeDark}
Ruta MILC: así vamos a aprender
\end{titlebox}
\begin{tabularx}{\textwidth}{>{\centering\arraybackslash}X>{\centering\arraybackslash}X>{\centering\arraybackslash}X>{\centering\arraybackslash}X}
\phasecard{1}{milcVerde}{milcGris}{Escucha\\[-1mm]\small Observo, escucho y pregunto.} &
\phasecard{2}{milcTurquesa}{milcGris}{Sistematización\\[-1mm]\small Pensamiento computacional.} &
\phasecard{3}{milcMagenta}{milcGris}{Praxis\\[-1mm]\small Construyo y aplico.} &
\phasecard{4}{milcVino}{milcGris}{Evaluación\\[-1mm]\small Reflexión liberadora.}
\end{tabularx}


\puente{Antes de teorizar, vas a hacer un \emph{inventario completo}: lista lo que tienes producido (todos los capítulos en V3, portada definitiva, mapa estructural, ficha de derechos, declaración de IA, archivo de diagramación) y lo que aún falta (capítulos pendientes, créditos, carta del editor). Ese inventario es el plan de hoy.}

\begin{titlebox}{milcVerde}
Fase 1 · Escucha
\end{titlebox}

\begin{softbox}{milcVerde}{milcGris}{Escena para escuchar}
\textbf{Actividad 1 · ANALIZA --- Inventario completo de lo producido} (15 min · individual). Haz lista de lo que tienes y de lo que falta para el libro de 80 páginas. Reglas: (1) \textbf{Tengo}: cada capítulo en V3 que esté listo (con apertura/cierre escritos a mano), la portada definitiva con composición, la ficha de derechos, el archivo de diagramación con plantilla aplicada. (2) \textbf{Me falta}: capítulos sin V3, página de créditos, tabla de contenido aplicada al libro completo, carta del editor, revisión final. Para cada item de la columna \emph{me falta}, anota cuánto tiempo necesitas para terminarlo.
\end{softbox}

\checkbox Listé todo lo que tengo producido.\\
\checkbox Listé todo lo que aún falta.\\
\checkbox Estimé tiempo necesario para cada pendiente.

\begin{infoband}{milcVerde}


\textbf{Lo que va al cuaderno (Actividad 1 · ANALIZA).} Título: «Actividad 1 --- Inventario de producción». \textbf{Formato:} 2 listas (tengo, me falta) + cronograma para cerrar pendientes esta semana. \textbf{Sabes que terminaste cuando} tienes claro qué hacer en qué orden.
\end{infoband}

\puente{Ya tienes el inventario. Ahora vas a entender la \textbf{anatomía del libro compilado profesional}: los 6 componentes obligatorios, el checklist de 12 puntos, los 5 errores típicos del editor que se apura en la producción final.}

\begin{titlebox}{milcTurquesa}
Fase 2 · Sistematización con pensamiento computacional
\end{titlebox}

\begin{softbox}{milcTurquesa}{milcGris}{Cuatro pilares aplicados al tema}
Tu inventario es el plan de la semana. Ahora necesitas la \textbf{anatomía profesional} del libro compilado para integrar todo con criterio.
\end{softbox}

\small
\begin{tabularx}{\textwidth}{>{\bfseries\RaggedRight\arraybackslash}p{3.6cm}Y}
\rowcolor{milcTurquesa!90}\color{white}Pilar & \color{white}\bfseries Aplicación al tema\\
1. Descomposición & Los \textbf{6 componentes obligatorios} del libro compilado se descomponen así: (1) \textbf{Portada}: imagen + título + autor. (2) \textbf{Página de créditos}: año de publicación, autor con nombre completo, licencia elegida (sesión 6), declaración de uso de IA con modelos y porcentaje aproximado, contacto del autor. (3) \textbf{Tabla de contenido} autogenerada con páginas correctas. (4) \textbf{8-12 capítulos diagramados}: cada uno con su título de estilo, texto V3, numeración correcta. (5) \textbf{Carta del editor}: 1 página, firmada con nombre y fecha, declarando proceso del libro y decisiones críticas. (6) \textbf{Contraportada} opcional con sinopsis de 200 palabras.\\
2. Reconocimiento de patrones & El \textbf{checklist de 12 puntos} de revisión final cubre 4 dimensiones: (a) \textbf{Contenido}: 8+ capítulos completos, V3 con voz propia, sin clichés evidentes. (b) \textbf{Forma}: tipografía coherente, jerarquía visual aplicada, numeración funcionando, márgenes respirables. (c) \textbf{Derechos}: créditos completos, licencia declarada, uso de IA explícito, fuentes citadas. (d) \textbf{Voz propia}: V3 con escritura humana en puntos críticos, sin que el libro suene a chatbot, carta del editor firmada.\\
3. Abstracción & La \textbf{carta del editor} es pieza profesional y emocional al mismo tiempo. Tiene 4 párrafos: (a) \textbf{Qué}: en 3-4 líneas, qué es este libro y para quién. (b) \textbf{Cómo}: el proceso seguido (concepción, prompts, iteración, voz propia, derechos). (c) \textbf{Declaración de IA}: qué modelos usaste, en qué partes, qué porcentaje aproximado de texto IA, qué porcentaje humano. (d) \textbf{Compromiso}: lo que asumes como autor responsable. Cierra con nombre completo, fecha, ciudad. Es tu firma profesional.\\
4. Algoritmo conceptual & Algoritmo: (1) cerrar capítulos pendientes en V3. (2) integrar todos los capítulos en el archivo de diagramación. (3) generar tabla de contenido autogenerada. (4) agregar portada como página 1. (5) agregar página de créditos como página 2. (6) escribir carta del editor. (7) pasar checklist de 12 puntos. (8) corregir lo que falle. (9) exportar PDF final. (10) verificar PDF abierto en otra computadora.\\
\end{tabularx}
\normalsize

\begin{drawbox}{milcTurquesa}{Anatomía del libro compilado}
\textbf{Portada:} imagen + título + autor. \\[1mm] \textbf{Créditos:} año + licencia + IA + contacto. \\[1mm] \textbf{TOC:} autogenerada. \\[1mm] \textbf{Capítulos:} 8-12 diagramados. \\[1mm] \textbf{Carta editor:} 1 página firmada. \\[1mm] \textbf{Contraportada:} opcional.
\end{drawbox}

\begin{softbox}{milcMostaza}{milcGris}{Errores comunes y su corrección}
(1) \textbf{Saltarse el checklist}: declarar como final lo que no se ha revisado. (2) \textbf{Tabla de contenido manual}: si añades una página, todo se rompe. (3) \textbf{Sin página de créditos}: el libro no declara derechos ni autor. (4) \textbf{Carta del editor genérica}: \emph{``espero les guste''} no es carta editorial. (5) \textbf{No verificar PDF en otra máquina}: a veces fuentes faltan en máquinas distintas y el libro se ve raro.
\end{softbox}

\begin{infoband}{milcTurquesa}


\textbf{Lo que va al cuaderno (Actividad 2 · EVALÚA).} Título: «Actividad 2 --- Anatomía del libro compilado». \textbf{Formato:} (a) los 6 componentes con su contenido; (b) checklist de 12 puntos imprimible; (c) plantilla de la carta del editor en 4 párrafos. \textbf{Sabes que terminaste cuando} podrías guiar a un compañero a compilar su libro con tu lista.
\end{infoband}

\puente{Tienes el método. Toca aplicarlo: integras componentes, pasas checklist, escribes carta del editor firmada, exportas PDF final, entregas bitácora del proceso completo del periodo.}

\begin{titlebox}{milcMagenta}
Fase 3 · Praxis con pensamiento computacional
\end{titlebox}

\begin{softbox}{milcMagenta}{milcGris}{Construyendo el producto con los cuatro pilares}
\textbf{Actividad 3 · CREA + EXPLICA --- Libro compilado + carta firmada + bitácora} (60 min en sesión + tiempo en casa para cerrar pendientes · individual). Hora de aplicar. Integras componentes, pasas checklist, escribes carta, exportas PDF final, llenas bitácora del proceso completo del periodo.
\end{softbox}

\small
\begin{tabularx}{\textwidth}{>{\bfseries\RaggedRight\arraybackslash}p{3.6cm}Y}
\rowcolor{milcMagenta!90}\color{white}Pilar & \color{white}\bfseries Aplicación al producto\\
Descomposición & Tu producción final se completa en 7 pasos: (1) cerrar capítulos pendientes en V3 (esta semana). (2) integrar todos los capítulos en el archivo de diagramación de la sesión 8. (3) agregar portada como página 1 y créditos como página 2. (4) generar tabla de contenido. (5) escribir y firmar carta del editor. (6) pasar checklist de 12 puntos uno a uno. (7) exportar PDF final y abrirlo en otra computadora o celular para verificar.\\
Patrones & Sigue el patrón \textbf{``revisar antes de firmar''}: nunca declares el libro final sin haber pasado los 12 puntos del checklist. La firma sin revisión es firma sin valor. La phronesis del editor consiste en \textbf{aceptar la disciplina del checklist aunque tome 1 hora adicional}.\\
Abstracción & La bitácora del proceso completo del periodo es trabajo aparte: 1 página resumiendo las 9 sesiones recorridas. Para cada sesión: 1 línea sobre qué hiciste, 1 línea sobre qué aprendiste. Es la memoria del oficio que te llevas al grado 11 y a la vida profesional. Documenta el viaje completo, no solo el producto final.\\
Algoritmo de armado & Algoritmo: (1) cerrar capítulos. (2) integrar diagramación. (3) agregar portada y créditos. (4) TOC. (5) escribir carta. (6) checklist. (7) corregir. (8) exportar PDF. (9) verificar en otra máquina. (10) llenar bitácora del periodo. (11) prepararse mentalmente para sustentación en S10.\\
\end{tabularx}
\normalsize

\begin{drawbox}{milcMagenta}{Checklist final de los 12 puntos}
\checkbox 8-12 capítulos completos en V3. \\[1mm] \checkbox Voz propia visible en aperturas y cierres. \\[1mm] \checkbox Sin clichés evidentes. \\[1mm] \checkbox Tipografía coherente en todo el libro. \\[1mm] \checkbox Jerarquía visual aplicada consistente. \\[1mm] \checkbox Numeración de páginas correcta. \\[1mm] \checkbox Márgenes e interlineado respirables. \\[1mm] \checkbox Página de créditos con licencia declarada. \\[1mm] \checkbox Uso de IA explícitamente declarado. \\[1mm] \checkbox Fuentes citadas si las hay. \\[1mm] \checkbox Carta del editor firmada con nombre y fecha. \\[1mm] \checkbox PDF abierto en otra máquina y verificado.
\end{drawbox}

\begin{softbox}{milcMagenta}{milcGris}{Plantilla de trabajo}
\textbf{Componentes.} \textit{Portada + créditos + TOC + 8-12 capítulos + carta + (contraportada).} \\[1mm] \textbf{Checklist.} \textit{12 puntos verificados.} \\[1mm] \textbf{Carta editor.} \textit{Qué + Cómo + IA + Compromiso, firmada.} \\[1mm] \textbf{Bitácora del periodo.} \textit{Resumen de las 9 sesiones.} \\[1mm] \textbf{PDF final.} \textit{Exportado y verificado.}
\end{softbox}

\begin{infoband}{milcMagenta}


\textbf{Lo que va al cuaderno (Actividad 3 · CREA + EXPLICA).} Título: «Actividad 3 --- Libro compilado + bitácora del periodo». \textbf{Formato:} (a) referencia al PDF final; (b) checklist marcado; (c) carta del editor firmada; (d) bitácora del periodo. \textbf{Sabes que terminaste cuando} podrías mostrarle el PDF a un familiar y firmarías con orgullo.
\end{infoband}

\puente{Lo que sigue resume \emph{qué} entregas hoy: PDF de 80+ páginas + checklist + carta firmada + bitácora del proceso. El criterio principal: que cualquiera abriendo tu PDF reconozca una pieza profesional, no documento escolar improvisado.}

\begin{titlebox}{milcMostaza}
Producto de la sesión
\end{titlebox}
\textbf{PDF 80+ páginas + checklist 12 puntos + carta editor firmada + bitácora}.
\begin{softbox}{milcMostaza}{milcGris}{Criterios de calidad}
(1) PDF compilado de 80+ páginas. (2) 6 componentes obligatorios presentes. (3) Checklist de 12 puntos pasado. (4) Carta del editor firmada con nombre y fecha. (5) Bitácora del periodo completa. (6) PDF verificado en otra máquina.
\end{softbox}

\puente{Antes de cerrar, mira la producción final desde las cinco dimensiones humanas. Firmar con honestidad lo que terminaste es honor del oficio; entregar sin revisar es desperdicio de las 8 sesiones anteriores.}

\begin{titlebox}{milcVino}
Fase 4 · Evaluación liberadora — cinco dimensiones
\end{titlebox}

\begin{softbox}{milcVino}{milcGris}{Sentido de la evaluación}
Evaluar no es solo revisar una nota. Es reconocer qué aprendí, qué cambió en mí, qué emociones manejé, cómo me relaciono con la ciudad y la comunidad, y qué le dejaré a quien viene detrás.
\end{softbox}

\textbf{Mi reflexión liberadora en cinco dimensiones.} Completa cada frase iniciada:

\small
\begin{tabularx}{\textwidth}{>{\bfseries\RaggedRight\arraybackslash}p{3.4cm}Y>{\RaggedRight\arraybackslash}p{4.6cm}}
\rowcolor{milcVino}\color{white}Dimensión & \color{white}\bfseries Mi respuesta (frame starter) & \color{white}\bfseries Algo que descubrí\\
1. Personal & Lo que más cambió en mí fue \linead{6.5cm} & \linead{4.2cm}\\
2. Emocional & La emoción más fuerte fue \linead{2.5cm} y la regulé \linead{3cm} & \linead{4.2cm}\\
3. Ciudadana & En democracia esto importa porque \linead{6cm} & \linead{4.2cm}\\
4. Local (Cartago) & En mi barrio sería útil para \linead{6.5cm} & \linead{4.2cm}\\
5. Intergeneracional & A mi abuela/abuelo le diría \linead{6.5cm} & \linead{4.2cm}\\
\end{tabularx}
\normalsize

\begin{infoband}{milcVino}
\textbf{Para la próxima sustentación.} Vuelve a esta tabla en seis meses. Verás que las respuestas cambian — no porque lo aprendido cambie, sino porque tú cambias. La evaluación liberadora es una conversación contigo a través del tiempo.
\end{infoband}


\puente{Y para terminar, tres voces sobre la firma. Dussel sobre la pieza terminada que devuelve a la comunidad, Marco Aurelio sobre la disciplina del checklist, Floridi sobre el libro firmado como acto profesional contemporáneo.}

\begin{titlebox}{milcGradeDark}
Triángulo de pensamiento · cierre filosófico
\end{titlebox}

\begin{softbox}{milcVino}{milcGris}{Dussel · lente del nosotros}
\textit{``El libro firmado por el editor humano en la era de IA es acto de dignidad cultural que la cultura generativa necesita.''}

Cuando firmas tu libro con nombre, fecha y declaración de uso de IA, estás haciendo lo que la cultura generativa contemporánea necesita: \textbf{seres humanos que asumen autoría con honestidad}. La filosofía de la liberación pide ese acto: no esconder lo hecho con IA, no fingir autoría pura, no diluir responsabilidad. Tu carta del editor firmada es modelo profesional contemporáneo que pocos autores aún practican.

\textbf{¿Mi carta del editor declara con honestidad el proceso, o intenta presentar el libro como obra puramente humana?}
\end{softbox}
\linead{\linewidth}\\[1mm]
\linead{\linewidth}

\begin{softbox}{milcMostaza}{milcGris}{Estoicismo · Marco Aurelio · cuidado interior}
\textit{``La disciplina del checklist antes de firmar es virtud profesional; el atajo de firmar sin revisar es vicio que se descubre con el tiempo.''}

Es tentador exportar el PDF y declararlo final apenas se cierra el último capítulo. La sabiduría estoica recomienda lo opuesto: \emph{pasa el checklist antes de la firma}. Esa hora de revisión adicional separa al editor profesional del aficionado. La phronesis del oficio honra a quien revisa con disciplina.

\textbf{¿Estoy pasando el checklist completo antes de declarar el libro final, o quiero terminar rápido y firmar?}
\end{softbox}
\linead{\linewidth}\\[1mm]
\linead{\linewidth}

\begin{softbox}{milcTurquesa}{milcGris}{Floridi · lente de la infoesfera}
\textit{``El libro firmado con declaración honesta de IA es el modelo editorial del siglo XXI; ocultar el uso de IA es modelo obsoleto.''}

En 2026 y los años siguientes, los libros producidos con asistencia de IA serán mayoría. La ética editorial del siglo XXI exige declaración honesta, no ocultamiento. Tu libro con carta del editor firmada se posiciona en el modelo profesional emergente. Aunque sea proyecto escolar, las prácticas son las mismas que la industria editorial profesional adoptará obligatoriamente en pocos años.

\textbf{¿Mi libro contribuye al modelo editorial transparente del siglo XXI, o sostiene el modelo opaco que ya está obsoleto?}
\end{softbox}
\linead{\linewidth}\\[1mm]
\linead{\linewidth}

\begin{titlebox}{milcVino}
Compromiso final
\end{titlebox}

\small
\begin{tabularx}{\textwidth}{>{\RaggedRight\arraybackslash}p{3.0cm}YYY}
\rowcolor{milcVino}\color{white}\bfseries Criterio & \color{white}\bfseries Lo logré & \color{white}\bfseries En proceso & \color{white}\bfseries Necesito apoyo\\
Comprensión del tema & Explico ideas centrales con mis palabras. & Reconozco ideas pero falta precisión. & Necesito repasar conceptos base.\\
Pensamiento computacional & Apliqué los 4 pilares al diseñar mi producto. & Apliqué algunos pilares. & Necesito releer la fase 2.\\
Aplicación y producto & Mi producto cumple los criterios y se puede revisar. & Producto funcional con ajustes pendientes. & Aún en construcción.\\
Reflexión 5 dimensiones & Respondí cada dimensión con honestidad. & Respondí algunas con detalle. & Mis respuestas son muy generales.\\
Triángulo de pensamiento & Conecté las 3 voces con mi trabajo real. & Conecté al menos una. & No relaciono las voces con mi proceso.\\
\end{tabularx}
\normalsize

\textbf{Hoy entendí que:}\\[1mm]
\linead{\linewidth}\\[1mm]
\linead{\linewidth}

\textbf{Mi compromiso para la próxima sesión es:}\\[1mm]
\linead{\linewidth}\\[1mm]
\linead{\linewidth}

\begin{softbox}{milcMostaza}{milcGris}{Cita de despedida · Marco Aurelio}
\textit{``El obstáculo en el camino se vuelve el camino. Lo que estorba al camino se vuelve camino.''}\\[1mm]
Cada error de hoy, bien observado, se convierte en pieza del camino que pisarás mañana.
\end{softbox}

\small
\begin{tabularx}{\textwidth}{>{\bfseries}p{3.6cm}Y}
\rowcolor{milcGradeDark!12}Nombre completo: & \linead{10cm}\\
Fecha: & \linead{4cm} \quad Curso/grupo: \linead{4cm}\\
Mi firma: & \linead{9cm}\\
\end{tabularx}
\normalsize

\begin{infoband}{milcGradeDark}
\textbf{Para la próxima sesión.} Llega con el PDF final del libro, la carta del editor firmada y la bitácora del periodo. En la sesión 10 vas a hacer la \textbf{sustentación pública + feria del libro escolar}: presentas tu libro al grupo, intercambian copias, leen fragmentos. Es el cierre ritual del proyecto editorial.
\end{infoband}

\end{document}
